% !TeX root = ../main.tex

\section{Introduction}
\label{sec:introduction}
One of the most startling features of textbook quantum theory is the fact that measurements do not reveal pre-existing variables, but rather \emph{create} them, in a sense.
There is no part of the theory which specifies, for example, where a particle is located.
The theory only provides probabilities that a certain measurement of, for example, position will produce a certain outcome~\cite[section 3]{sep-qm}.

In 1935, Albert Einstein, Boris Podolsky, and Nathan Rosen (EPR) published a paper in which they considered a special kind of quantum system, a pair of \emph{entangled} particles~\cite{Einstein:1935rr}.
They argued that the special correlations such a system produces could only be explained if measurements do, in fact, \emph{reveal} pre-existing variables.
That is, they claim that quantum theory does not describe reality completely, and should be replaced by a theory which includes these \emph{hidden} variables.

In his 1951 textbook, David Bohm presented a conceptually simpler version of the EPR argument~\cite[chapter 22]{Bohm:1951}.
This EPRB argument is formulated in terms of a pair of entangled spin-1/2 particles in the singlet state.
According to quantum theory, the measurements of the spin components along a specified axis will perfectly anti-correlate.

It should be noted that, as Don Howard documented in 1985, \blockcquote[172]{HOWARD1985171}{Einstein did not write the Einstein, Podolsky and Rosen paper, and he was unhappy with the way it turned out}.
Einstein demonstrated his argument in a 1948 paper, starting from a \emph{principle of locality}~\cite{Kiefer2022}.\footnote{
    Howard also identifies an assumption of \emph{separability}.
    See also the entry on \enquote{Holism and Nonseparability in Physics} in The Stanford Encyclopedia of Philosophy, \url{https://plato.stanford.edu/archives/spr2022/entries/physics-holism/}.
}
Thus, the EPR argument shows an implication from locality to the existence of hidden variables.

This line of argument was continued by John Bell in the 1964 derivation of his famous inequality~\cite{Bell:1964kc}.
He showed that, given certain auxiliary assumptions, the existence of these local hidden variables cannot be reconciled with the predictions of quantum theory.
Specifically, measurements of the spin components along axes at arbitrary angles will prove decisive.

\emph{Bell's theorem}, presented in the same 1964 paper, combines the EPRB argument (from locality to pre-existing values) and his inequality (from pre-existing values to a violation of the predictions of quantum theory).
It follows that any theory which reproduces the predictions of quantum theory cannot always be local.

We will present an inequality first derived by Clauser, Horne, Shimony, and Holt~\cite{Clauser:1969ny}, which was subsequently refined by Bell~\cite{Bell:1971}.
In this version, the assumption of perfect correlations was relaxed, thereby making the system more experimentally accessible.
Furthermore, this inequality places restrictions on any theory which \emph{approximately} reproduces the predictions of quantum theory.

In the last two sections, we will examine the necessary assumptions and review experimental tests of the situations at hand.
The central assumption is a \emph{principle of locality}, but there are also important auxiliary assumptions.
The experimental tests faced certain \enquote{loopholes}, which made a definitive interpretation difficult.
The following presentation will be mainly influenced by a 1980 paper by Bell, and the entry on Bell's theorem in \emph{The Stanford Encyclopedia of Philosophy}~\cite{Bell:1980wg,sep-bell-theorem}.